% ============================================================
%  Eat Consciously — A practical guide to food additives
%  faustobe.it | 2026 | CC BY-NC 4.0
% ============================================================
\documentclass[a5paper,12pt,oneside,openany]{book}

\usepackage[T1]{fontenc}
\usepackage[utf8]{inputenc}
\usepackage{ebgaramond}
\usepackage[a5paper, top=2cm, bottom=3cm, inner=2.2cm, outer=3cm]{geometry}
\usepackage[table]{xcolor}
\usepackage[english]{babel}

% Colours
\definecolor{verdescuro}{HTML}{2d6a4f}
\definecolor{verdecharo}{HTML}{52b788}
\definecolor{grigioscuro}{HTML}{2b2d42}
\definecolor{warnyellow}{HTML}{d4a017}
\definecolor{warnred}{HTML}{c1121f}
\definecolor{infoBlue}{HTML}{3a86ff}

\usepackage{tcolorbox}
\tcbuselibrary{skins,breakable}

\usepackage{booktabs}
\usepackage{tabularx}
\newcolumntype{L}{>{\raggedright\arraybackslash}X}

\usepackage{hyperref}
\hypersetup{
  colorlinks=true,
  linkcolor=verdescuro,
  urlcolor=verdescuro,
  pdftitle={Eat Consciously},
  pdfauthor={faustobe},
}

\usepackage{qrcode}
\usepackage{fancyhdr}
\usepackage{titlesec}
\usepackage{microtype}

% Coloured boxes
\tcbset{
  boxbase/.style={
    breakable, enhanced, arc=4pt, boxrule=0.8pt,
    fonttitle=\bfseries\small,
    before upper={\parskip=4pt},
  }
}

\newcommand{\boxgreen}[2]{%
  \begin{tcolorbox}[boxbase,
    colback=verdescuro!8!white, colframe=verdescuro, title={#1}]
  #2\end{tcolorbox}}

\newcommand{\boxyellow}[2]{%
  \begin{tcolorbox}[boxbase,
    colback=warnyellow!12!white, colframe=warnyellow!80!black, title={#1}]
  #2\end{tcolorbox}}

\newcommand{\boxred}[2]{%
  \begin{tcolorbox}[boxbase,
    colback=warnred!8!white, colframe=warnred, title={#1}]
  #2\end{tcolorbox}}

\newcommand{\boxinfo}[2]{%
  \begin{tcolorbox}[boxbase,
    colback=infoBlue!8!white, colframe=infoBlue, title={#1}]
  #2\end{tcolorbox}}

% Page style
\pagestyle{fancy}
\fancyhf{}
\fancyfoot[C]{\footnotesize\textcolor{grigioscuro}{\textit{Eat Consciously\ \textbar\ faustobe.it}}}
\fancyfoot[R]{\footnotesize\thepage}
\renewcommand{\headrulewidth}{0pt}
\renewcommand{\footrulewidth}{0.4pt}
\fancypagestyle{plain}{%
  \fancyhf{}%
  \fancyfoot[C]{\footnotesize\textcolor{grigioscuro}{\textit{Eat Consciously\ \textbar\ faustobe.it}}}%
  \fancyfoot[R]{\footnotesize\thepage}%
  \renewcommand{\headrulewidth}{0pt}%
  \renewcommand{\footrulewidth}{0.4pt}%
}

% Chapter title style
\titleformat{\chapter}[display]
  {\normalfont\LARGE\bfseries\color{verdescuro}}
  {\chaptertitlename\ \thechapter}{10pt}{\Huge}
\titlespacing*{\chapter}{0pt}{-20pt}{30pt}

% ============================================================
\begin{document}

% -------- COVER --------
\begin{titlepage}
\pagecolor{verdescuro}
\color{white}
\vspace*{2.5cm}
\begin{center}
  {\fontsize{38}{44}\selectfont\bfseries Eat Consciously}\\[1.2cm]
  \textcolor{verdecharo}{\rule{6cm}{0.6pt}}\\[1.2cm]
  {\large\itshape A practical guide to food additives}\\[0.6cm]
  {\normalsize With E-Codes Reader for Android}\\[3cm]
  {\normalsize\textcolor{verdecharo}{faustobe.it}}
\end{center}
\vfill
\begin{center}
  {\small Version 1.0\ \textbullet\ 2026\ \textbullet\ CC BY-NC 4.0}
\end{center}
\end{titlepage}
\nopagecolor

\frontmatter
\tableofcontents
\mainmatter

% ============================================================
\chapter{How to use the app for conscious shopping}

Learn to recognise the most processed foods and choose genuinely simple,
natural ones in just a few minutes --- without becoming an additive expert.

\section{Before scanning, do three quick checks}

Even before opening the app, look at the packaging. Front-of-pack labels
are often designed to make products appear healthier than they really are.

\begin{itemize}
  \item \textbf{Look at the front of the pack.} Watch out for claims like
    ``sugar-free but with sweeteners'', ``fat-free but with emulsifiers'',
    ``extra flavour'' or ``flavoured''. These often signal that the product
    has been heavily modified.
  \item \textbf{Open the app and scan the product.} Use the camera to scan
    the barcode or upload a photo of the label. The app will immediately
    show you the ingredient list and any additives present.
  \item \textbf{Check three key points.} How many additives appear
    (especially E codes)? How long and complex is the ingredient list?
    How much sugar, salt and fat is declared?
\end{itemize}

If you answer ``many'' or ``long/complicated'' to more than one question,
you are probably looking at a highly processed food.

\section{Three rules to quickly spot an ``over-processed'' food}

\textbf{Fewer ingredients = better.} If the list is short and contains
familiar names (e.g.\ ``flour, water, salt, oil''), the product is minimally
processed and generally more trustworthy.

\textbf{Many ``unpronounceable'' ingredients.} Terms like ``modified
starches'', ``isolated proteins'', ``mono- and diglycerides of fatty acids'',
``flavour enhancers'' or ``complex flavourings'' are a clear sign of a
heavily processed food.

\textbf{Lots of E codes and additives.} Not all E numbers are harmful, but
their sheer number --- colourings, preservatives, stabilisers,
emulsifiers --- signals a high degree of industrial processing.

\section{How to read the app's result}

Each product receives a quick rating --- think of these colours as a food
traffic light.

\boxgreen{Minimally processed --- go ahead}{Few ingredients, few or no
E additives. Moderate sugar and salt. A great choice for everyday
shopping.}

\boxyellow{Moderately processed --- use in moderation}{Some E numbers,
natural flavourings or added sugars. Fine occasionally, not as a staple
of every meal.}

\boxred{Ultra-processed --- limit or avoid}{Many additives, added sugars,
complex flavourings and industrially processed fats. Best kept as a rare
exception, especially for children.}

\section{How to use the app practically while shopping}

\textbf{1.~Start with the chilled and breakfast aisles.} Scan yoghurts,
bread, milk and drinks. Look for products with few ingredients, few
E numbers and simple sugars (fruit, milk, cane sugar) rather than syrups
and sweeteners.

\textbf{2.~Then move on to sweets and snacks.} Check biscuits, salty
snacks and packaged desserts. If the list is very long and full of
E codes, the app will flag red. In those cases, always choose the simpler
alternative (fruit, wholegrain bread, nuts).

\textbf{3.~Use the app to compare similar products.} Find two products of
the same type and compare their app profiles. Often the one with fewer
ingredients is only slightly less eye-catching on the shelf but far more
natural.

\textbf{4.~Keep ``reds'' as exceptions, not the rule.} There is no need
to eliminate everything ultra-processed --- just reduce it. Set yourself
a limit --- for example, no more than 1--2 ``red'' products per week ---
while keeping the rest of your shop focused on simple foods.

\section{Summary table}

{\small
\begin{tabularx}{\textwidth}{|L|L|}
\hline
\rowcolor{verdescuro}\textcolor{white}{\textbf{App result}}
  & \textcolor{white}{\textbf{What to do in practice}} \\
\hline
Few ingredients, few/no E numbers
  & Use as everyday staple: plain yoghurt, wholegrain bread, tomato passata \\
\hline
Some E numbers, moderate sugar
  & Occasionally, not as a main food \\
\hline
Many E numbers, added sugars, complex flavourings
  & Limit or avoid: keep it as a rare exception \\
\hline
\end{tabularx}
}

\section{Why learning to read additives really makes a difference}

Recognising highly processed foods is not ``additive phobia'' --- it is
\textbf{paying attention to the quality of what we eat}. Many studies
suggest that reducing ultra-processed foods helps with long-term weight
management, cardiovascular health and metabolic stability.

The app gives you a practical tool: you do not need to memorise every
E code, just learn to spot the signs of excessive processing and trust
the quick rating.

% ============================================================
\chapter{How to read the nutrition label}

Ingredients, calories, hidden salt: learn to decode a label in under a
minute and stop buying products that look healthy but are not.

\section{The ingredients list: where to look first}

By law, ingredients are listed in \textbf{descending order by weight}:
what appears first is present in the greatest amount. This is the most
important starting point on any label.

\begin{itemize}
  \item \textbf{The first ingredient matters most.} If the first ingredient
    is ``sugar'', ``glucose syrup'' or ``palm oil'', the product is built
    around that element.
  \item \textbf{List length = degree of processing.} Homemade bread: flour,
    water, salt, yeast --- 4 ingredients. Industrial supermarket bread:
    can have 15 or more.
  \item \textbf{Look for names that are hard to pronounce.} Terms like
    ``modified starch'', ``mono- and diglycerides of fatty acids'',
    ``carboxymethyl cellulose'' or ``hydrogenated'' indicate industrially
    derived ingredients.
\end{itemize}

\boxinfo{Practical rule}{If you could not find all the ingredients at a
regular grocery store, the product is probably highly processed.}

\section{The nutrition table: what to look at (and what to ignore)}

\begin{description}
  \item[Sugars] Below 5\,g per 100\,g is low; between 5 and 22.5\,g is
    medium; above 22.5\,g is high. For drinks, the ``low'' threshold is
    2.5\,g/100\,ml.
  \item[Salt] Below 0.3\,g/100\,g is low; above 1.5\,g/100\,g is high.
    Sodium on the label must be multiplied by 2.5 to get the salt value.
  \item[Saturated fat] Below 1.5\,g/100\,g is acceptable; above
    5\,g/100\,g is considered high.
\end{description}

\section{Servings vs 100\,g: the trick that confuses everyone}

Nutritional values are often shown both per 100\,g and ``per serving''.
Companies tend to highlight the per-serving values, which look lower ---
but serving sizes are frequently underestimated.

\boxyellow{Example}{If the pack declares 10\,g of sugar per serving (30\,g)
and you eat a 60\,g bowl, you are actually taking in 20\,g of sugar ---
nearly half the recommended daily limit.}

\section{Hidden salt: where it really lurks}

Salt is not only found in savoury products. In many apparently sweet foods
--- biscuits, cereals, bread --- the salt content is surprisingly high.

\begin{itemize}
  \item \textbf{Bread and baked goods:} two slices of industrial bread can
    contain 0.5--0.8\,g of salt.
  \item \textbf{Breakfast cereals:} even ``healthy'' wholegrain varieties
    can contain 0.8--1.2\,g of salt per 100\,g.
  \item \textbf{Sauces and condiments:} ready-made sauces can contain
    between 1 and 3\,g of salt per 100\,g.
\end{itemize}

\boxinfo{Important note}{On the label you may find ``sodium'' instead of
``salt''. To convert: \textbf{salt = sodium $\times$ 2.5}.}

\section{Nutrition claims: what they really mean}

Claims like ``fat-free'', ``light'', ``high in fibre'' or ``0\% sugar'' are
regulated by law but can still be misleading.

\begin{itemize}
  \item \textbf{``No added sugars''} does not mean sugar-free: it may
    still contain naturally occurring sugars (fructose, lactose) or
    artificial sweeteners.
  \item \textbf{``Light'' or ``reduced''} means reduced by 30\% compared
    to the reference product. Often the fat removed is replaced with sugar
    or thickeners.
  \item \textbf{``High in fibre''} requires at least 6\,g/100\,g but says
    nothing about the other ingredients: a ``high-fibre'' cookie can still
    contain a lot of sugar.
\end{itemize}

\section{Quick checklist to use at the supermarket}

\begin{enumerate}
  \item Look at the first ingredient: is it a real food?
  \item Count the ingredients: fewer than 5 is great; more than 10--12,
    take a closer look.
  \item Check the E codes: few or none is a good sign.
  \item Check salt and sugars per 100\,g: salt above 1.5\,g or sugars
    above 22.5\,g are high.
  \item Ignore the big claims on the front: always trust the ingredients
    list, not the promotional text.
\end{enumerate}

% ============================================================
\chapter{Additives vs ultra-processed foods: what's the difference?}

Having many E codes is not the same as being ultra-processed.
Understanding the difference helps you make more complete and accurate
assessments.

\section{What food additives are}

Food additives are substances intentionally added to food to perform a
technological function: preserving, colouring, thickening, emulsifying,
stabilising or enhancing flavour. In Europe, every approved additive
receives an \textbf{E code} and must pass a safety assessment by EFSA.

\begin{itemize}
  \item \textbf{Not all additives are synthetic.} Citric acid (E330),
    curcumin (E100), soy lecithin (E322) and ascorbic acid (E300,
    vitamin~C) are naturally derived additives.
  \item \textbf{Not all additives are dangerous.} The vast majority of
    approved E codes are considered safe at permitted usage levels.
  \item \textbf{Some additives deserve attention.} Azo dyes (E102, E110,
    E122, E124, E129), nitrates (E249--E252), and intense sweeteners
    (E950--E962) have more consistent risk evidence, especially for
    children.
\end{itemize}

\section{What ultra-processed foods are}

The term ``ultra-processed'' comes from the \textbf{NOVA classification},
developed at the University of S\~{a}o Paulo. It is not about nutrients,
but about the \emph{industrial process} that produced the food.

A food is ultra-processed (NOVA 4) when it contains ingredients you
would not find in a home kitchen: isolated proteins, modified starches,
glucose-fructose syrups, ``nature-identical'' flavourings, emulsifiers,
colourings, stabilisers --- often in complex combinations.

\section{When they overlap --- and when they don't}

\boxgreen{Additives but NOT ultra-processed}{Cheese with preservative
E252: processed but not NOVA~4. Wine with sulphites (E220): natural
additive, traditional product. Oil with vitamin~E (E306): minimal
processing.}

\boxred{Ultra-processed with FEW E additives}{Toast bread with modified
starch and flavourings: not many E codes, but still NOVA~4. Protein bars
with isolated proteins and rice syrup: few codes, but heavily
transformed.}

\boxyellow{Ultra-processed AND many additives (the worst case)}{Snack
cakes: modified starches + colourings + emulsifiers + flavourings +
sweeteners. Industrial sauces: preservatives + stabilisers + flavour
enhancers + colourings.}

\section{The double problem: why both matter}

\begin{itemize}
  \item \textbf{Risks linked to specific additives:} azo dyes and
    hyperactivity in children (EFSA 2007); nitrates in processed meat
    and carcinogenic nitrosamines; intense sweeteners and possible gut
    microbiome disruption.
  \item \textbf{Risks linked to ultra-processing itself:} large-scale
    studies (NutriNet-Santé, UK~Biobank) associate high NOVA~4 consumption
    with obesity, type~2 diabetes, cardiovascular disease and some cancers
    --- independently of nutritional quality.
  \item \textbf{The combination effect:} additive studies almost always
    test one substance at a time. In real life we eat combinations of
    10--20 different additives in the same meal.
\end{itemize}

\section{How to spot them in practice: the double read}

\begin{enumerate}
  \item Look for ``lab-derived'' ingredients: modified starches, isolated
    proteins, glucose-fructose syrups, ``nature-identical'' flavourings.
  \item Count the E codes and identify critical categories: a natural
    antioxidant (E306) is very different from an azo dye (E102).
  \item Consider the meal context: a balanced meal with vegetables and
    legumes can better accommodate the occasional more processed product.
  \item Use the app as a guide, not the final judge.
\end{enumerate}

\boxinfo{Summary in one sentence}{A product with one natural E code is
not a problem. A product with 15 industrial ingredients and no E codes
at all is still ultra-processed. Look at both things.}

% ============================================================
\chapter{NOVA classification: what it is and how to use it}

The NOVA system classifies foods by degree of industrial processing ---
not by calories or nutrients. Understanding the four groups helps you
make better choices every day.

\section{What the NOVA classification is}

NOVA is a food classification system developed by the research group of
Professor Carlos Monteiro at the University of S\~{a}o Paulo (Brazil)
and now adopted by the WHO, FAO and numerous global epidemiological
studies.

Unlike nutrient-based systems (such as Nutri-Score), \textbf{NOVA
classifies foods according to their degree of industrial processing}:
not ``how many calories does it have'', but ``how much industrial
transformation went into making it''.

\section{Group 1 --- Unprocessed or minimally processed foods}

Foods of plant or animal origin that have not undergone any significant
industrial transformation, or only minimal physical processes (drying,
milling, chilling, pasteurisation).

\textbf{Examples:} fresh or frozen fruit and vegetables, fresh or frozen
meat and fish without additives, eggs, dried pulses, whole-grain cereals
(rice, oats, spelt), fresh pasteurised milk, plain yoghurt without
additives.

Group 1 foods should form the base of the diet.

\section{Group 2 --- Processed culinary ingredients}

Substances extracted from Group~1 foods or from nature, used to cook,
season and prepare food. Not eaten on their own.

\textbf{Examples:} olive oil and other vegetable oils, butter, vinegar,
salt, sugar, honey, plain flour, pure cornstarch, maple syrup. Use them
in moderation.

\section{Group 3 --- Processed foods}

Products made by adding Group~2 ingredients to Group~1 foods, using
processes such as salt-preservation, smoking, fermentation or ageing.
Recognisable because they typically contain just 2--3 ingredients.

\textbf{Examples:} canned vegetables and pulses, quality cured ham,
traditional cheeses (parmesan, pecorino), artisan bread, wine.

\section{Group 4 --- Ultra-processed: the group to limit}

Industrial formulations that contain no whole food ingredients, or use
them only minimally. The key signal: they contain ingredients
\textbf{you would not find in a home kitchen}.

\textbf{Typical Group 4 ingredients:} hydrolysed proteins, modified
starches, glucose-fructose syrup, hydrogenated oils, ``nature-identical''
flavourings, artificial colourings, intense sweeteners (aspartame,
saccharin, sucralose), emulsifiers (E471, E472).

\boxred{Key data}{In Italy around 14--18\% of daily calories come from
ultra-processed foods. In the UK and USA the figure exceeds 50--60\%.
Even consuming more than 20\% of daily calories from NOVA~4 foods is
associated with health risks.}

\section{NOVA in everyday shopping}

{\small
\begin{tabularx}{\textwidth}{|L|L|L|L|}
\hline
\rowcolor{verdescuro}
\textcolor{white}{\textbf{Group}}
  & \textcolor{white}{\textbf{Definition}}
  & \textcolor{white}{\textbf{Examples}}
  & \textcolor{white}{\textbf{In diet}} \\
\hline
NOVA~1 & Unprocessed/min.\ processed & Fruit, veg, eggs, fresh meat & Foundation \\
\hline
NOVA~2 & Culinary ingredients & Oil, salt, sugar, flour & In moderation \\
\hline
NOVA~3 & Traditionally processed & Preserves, cheeses, artisan bread & Acceptable \\
\hline
NOVA~4 & Ultra-processed & Snack cakes, hot dogs, sugary cereals & Limit \\
\hline
\end{tabularx}
}

% ============================================================
\chapter{Ingredients to watch out for on labels}

Hidden sugars, dangerous fats, controversial additives and excess sodium:
learn to spot them at a glance.

\section{Sugars: the 10 hidden names on labels}

Sugar does not always go by the name ``sugar''. Companies use dozens of
variations to spread it across multiple entries in the ingredients list.

\begin{itemize}
  \item \textbf{Glucose-fructose syrup} (the worst: metabolised
    differently from table sugar)
  \item \textbf{Maltodextrins} (starch broken down almost completely)
  \item \textbf{Dextrose}, \textbf{Fructose}
  \item \textbf{Concentrated cane juice}, \textbf{Corn syrup}
  \item \textbf{Rice syrup}, \textbf{Dehydrated honey}
  \item \textbf{Coconut sugar}, \textbf{Agave nectar}
\end{itemize}

The WHO recommends keeping free sugars below \textbf{10\% of daily
calories} --- about 50\,g per day for an adult on 2\,000\,kcal.
Preferably below 5\% (25\,g).

\section{Fats to avoid: hydrogenated and trans fats}

Trans fatty acids, produced by the industrial hydrogenation of vegetable
oils, are associated with increased cardiovascular risk.

\boxred{Hydrogenated and partially hydrogenated fats}{Look on the label
for: ``partially hydrogenated vegetable oil'', ``hydrogenated vegetable
fat'', ``shortening''. In Europe they have been banned above 2\,g/100\,g
of fat since 2021. Avoid them regardless.}

\boxyellow{Palm oil: not trans, but\ldots}{During refining at high
temperatures, potentially carcinogenic compounds form (3-MCPD, glycidol).
It is not banned, but worth limiting.}

\section{Controversial additives: the E codes to know}

\begin{itemize}
  \item \textbf{Azo dyes --- maximum caution with children:} E102
    (Tartrazine), E104, E110, E122, E124, E129. EFSA study 2007: link to
    hyperactivity in children. In Europe a mandatory label warning is
    required: ``may have an adverse effect on activity and attention in
    children''.
  \item \textbf{Nitrates and nitrites --- processed meat:} E249--E252.
    In the presence of animal proteins they form nitrosamines, classified
    as probable carcinogens (IARC Group~2A).
  \item \textbf{Intense sweeteners:} E950 (Acesulfame K), E951
    (Aspartame), E952 (Cyclamate), E954 (Saccharin), E955 (Sucralose).
    The WHO (2023) expressed caution about chronic use.
\end{itemize}

\section{Hidden sodium: beyond the salt shaker}

Excess sodium raises blood pressure. Around 75\% of sodium in a Western
diet does not come from what you add during cooking --- it is already in
the packaged product.

\textbf{Main sources:} bread, cheeses, cured meats, ready meals, jarred
sauces, stock cubes, soy sauce, ketchup, salted snacks, industrial
breakfast cereals.

Additives such as monosodium glutamate (E621) and sodium benzoate (E211)
also contribute to total sodium intake.

\section{The rule of 5: the fastest filter}

\begin{enumerate}
  \item More than 5 ingredients? Start reading carefully.
  \item Sugars above 5\,g/100\,g? Check where they come from.
  \item Salt above 0.5\,g/100\,g in a sweet product? A sign of
    ultra-processing.
  \item More than 5 E additives? The product is probably ultra-processed.
  \item Do not recognise 5 or more ingredients? Put it back on the shelf.
\end{enumerate}

% ============================================================
\chapter{Recommended foods}

Substitute, don't eliminate: simple, accessible and often cheaper
alternatives to the most common ultra-processed foods.

\section{The basic principle: substitute, don't eliminate}

This is not about strict diets or feeling guilty. It is about knowing the
alternatives --- and choosing them when they are as available and practical
as the product they replace.

\textbf{The 80/20 rule works:} if 80\% of what you eat is Group~1--2
(NOVA), the 20\% of more processed products makes no real difference.
Dietary health is measured on a weekly average, not on a single meal.

\section{Grains and carbohydrates: the best choices}

\textbf{Recommended grains:} rolled oats (no added ingredients), brown
rice, spelt grain, pearl barley, buckwheat, wholemeal pasta made only
from wholemeal semolina, bread made only from wholemeal flour (maximum
4--5 ingredients).

\textbf{Limit these:} sugary breakfast cereals, sandwich bread with
modified starches and preservatives, crackers with palm oil, pre-cooked
rice and pasta in pouches with additives.

\boxinfo{Quick tip for breakfast cereals}{If sugars are above 10\,g/100\,g,
you are eating a dessert product --- not a healthy breakfast.}

\section{Protein: the best sources}

\begin{itemize}
  \item \textbf{Fish and seafood:} fresh or frozen fish (without batter),
    sardines and mackerel in olive oil, tuna in water, salt-packed
    anchovies (rinsed).
  \item \textbf{Eggs and dairy:} whole eggs, plain Greek yoghurt, ricotta,
    cottage cheese, aged cheeses in moderation.
  \item \textbf{Legumes:} lentils, chickpeas, beans, peas. Canned in
    water, they are already cooked and can be added to any meal in 30
    seconds.
\end{itemize}

\section{Good fats: the easiest ones to add}

\textbf{Recommended sources:} extra-virgin olive oil, raw linseed oil,
walnuts, almonds, hazelnuts, flaxseeds, chia seeds, avocado, oily fish
(salmon, mackerel, sardines, tuna), whole eggs.

\textbf{In moderation:} butter, coconut oil, whole milk.

\textbf{Fats to avoid:} margarines with partially hydrogenated fats,
vegetable shortening, refined palm oil in industrial baked goods, vegetable
cream with emulsifiers.

\section{A sample shopping list}

{\small
\begin{tabularx}{\textwidth}{|L|L|}
\hline
\rowcolor{verdescuro}\textcolor{white}{\textbf{Category}}
  & \textcolor{white}{\textbf{Recommended products}} \\
\hline
Grains & Rolled oats, brown rice, spelt, wholemeal pasta, wholemeal artisan bread \\
\hline
Legumes & Lentils, chickpeas, beans, peas (dried or canned in water) \\
\hline
Animal protein & Eggs, sardines in olive oil, tuna in water, plain Greek yoghurt, ricotta \\
\hline
Vegetables & Seasonal, fresh or frozen without additives \\
\hline
Fats & Extra-virgin olive oil, walnuts, almonds, flaxseeds \\
\hline
Fruit & Whole seasonal fruit, frozen berries (no added sugar) \\
\hline
Condiments & Apple cider vinegar, minimally processed soy sauce, herbs and spices \\
\hline
\end{tabularx}
}

% ============================================================
\chapter{Why ultra-processed foods are harmful}

The scientific evidence on health risks linked to habitual consumption of
ultra-processed foods: obesity, inflammation, microbiota, chronic disease.

\section{Obesity and metabolic disruption}

Ultra-processed foods do not cause weight gain simply because they are
high in calories. The mechanism is more insidious: they are engineered to
be \textbf{hyper-palatable} and actively interfere with the body's natural
satiety signals.

\boxinfo{Kevin Hall's study (NIH, 2019)}{In a controlled clinical trial,
participants exposed to an ultra-processed diet consumed an average of
\textbf{500 extra calories per day} compared to those eating minimally
processed foods --- with equal freedom of choice. In just two weeks,
the UPF group had gained nearly one kilogram.}

\begin{itemize}
  \item \textbf{High energy density:} UPFs contain on average
    2.15\,kcal per gram --- almost double that of fresh foods.
  \item \textbf{Excessive stimulation of pleasure centres:} the industrial
    combination of sugar, salt, fat and flavourings activates the
    dopaminergic system. Eating continues even in the absence of real
    hunger.
  \item \textbf{Reduction of peptide YY:} UPFs lower levels of the gut
    hormone that signals fullness to the brain.
\end{itemize}

\boxred{+41\% risk of abdominal obesity}{People who follow a diet high in
UPFs have a 41\% higher risk of developing obesity compared to those who
consume few of them.}

\section{Inflammation and gut microbiota}

Frequent consumption of ultra-processed foods is associated with an
increase in inflammatory immune responses. Emulsifiers such as
carboxymethylcellulose (E466) and polysorbate~80 (E433) alter the
protective intestinal mucus even at doses permitted by regulations.

The ultra-processing method strips foods of fibres and bioactive plant
compounds, promoting \textbf{dysbiosis} --- a gut microbiota imbalance
associated with insulin resistance and chronic inflammation.

\section{A systemic risk: multiple mechanisms at once}

{\small
\begin{tabularx}{\textwidth}{|L|L|L|}
\hline
\rowcolor{verdescuro}
\textcolor{white}{\textbf{Mechanism}}
  & \textcolor{white}{\textbf{Effect}}
  & \textcolor{white}{\textbf{Consequence}} \\
\hline
Hyper-palatability & Excessive calorie intake & Obesity, type~2 diabetes \\
\hline
Fibre deficiency & Intestinal dysbiosis & Chronic inflammation, IBD \\
\hline
Additives (emulsifiers) & Intestinal permeability & Autoimmune diseases \\
\hline
Sugars and sweeteners & Insulin resistance & Type~2 diabetes \\
\hline
Replacement of fresh food & Micronutrient deficiency & Oxidative stress \\
\hline
\end{tabularx}
}

\boxinfo{Large epidemiological studies confirm this}{NutriNet-Sant\'{e}
(France, over 100,000 participants) and UK~Biobank (UK, over 500,000
participants) associate high UPF consumption with greater risks of cancer,
cardiovascular disease, depression and all-cause mortality ---
independently of the nutritional quality of individual products.}

\section{What to do in practice}

\begin{enumerate}
  \item \textbf{Identify your habitual UPFs:} sweetened cereals, packaged
    snacks, sliced bread with additives, ready meals, flavoured drinks.
  \item \textbf{Substitute one at a time:} rolled oats instead of sweet
    cereals; plain Greek yoghurt instead of flavoured yoghurt.
  \item \textbf{Increase fibre --- it is the priority:} legumes, vegetables,
    whole-grain cereals, whole fruit.
  \item \textbf{Use the app to identify hidden UPFs:} many products that
    appear healthy are actually NOVA~4.
  \item \textbf{Do not aim for perfection:} even a 20--30\% improvement
    in diet quality produces measurable health effects in the medium term.
\end{enumerate}

% ============================================================
\chapter{Food safety for children and the elderly}

Same dose, different effects: why the most vulnerable groups require
particular attention to additives and sodium.

\section{Why children and the elderly are more vulnerable}

The same dose of an additive that a healthy adult metabolises without
issue can have different effects on a young child or an elderly person.
This is not alarmism: it is physiology.

\begin{itemize}
  \item \textbf{Children: low body weight, immature systems.} The
    acceptable daily intake (ADI) is calculated in mg per kg of body
    weight. A 15\,kg child eating a 30\,g snack with colourings receives
    a relative dose 4--5 times higher than a 70\,kg adult. The liver and
    kidneys of young children also metabolise certain substances more
    slowly.
  \item \textbf{The elderly: reduced kidney and liver function.} With
    age, the capacity to eliminate certain substances decreases. Excess
    sodium has a greater impact on those with hypertension or kidney
    insufficiency.
  \item \textbf{Cumulative effects:} an ice cream with colourings once
    a month is not a problem. If children daily eat products containing
    E102, E110, E122 at breakfast, as a snack and at lunch --- the
    cumulative effect becomes significant.
\end{itemize}

\section{Critical additives for children}

\boxred{Azo dyes --- avoid with children}{E102 (Tartrazine), E104,
E110, E122, E124, E129. Mandatory European label warning: \emph{``may
have an adverse effect on activity and attention in children''}. Found in:
coloured sweets, orange/red fizzy drinks, jellies, industrial cakes,
coloured breakfast cereals.}

\boxyellow{Intense sweeteners --- not approved for young children}{E951
(Aspartame), E950 (Acesulfame K), E955 (Sucralose), E954 (Saccharin).
Not permitted in foods specifically for infants (EU Reg.~1333/2008).
Often found in ``0\% sugar'' yoghurts, ``light'' or ``zero'' drinks,
sugar-free sweets.}

\section{What to limit for the elderly}

\begin{itemize}
  \item \textbf{Sodium: the main risk.} Ready meals, cured meats, aged
    cheeses and industrial stock are the main sources. Target: less than
    5\,g of salt per day.
  \item \textbf{Phosphates: they interfere with calcium and bones.}
    E338--E341, E450--E452 are common in processed meat and processed
    cheese. High consumption interferes with calcium absorption.
  \item \textbf{Interactions with common medications:} vitamin~K in leafy
    greens with warfarin, grapefruit with statins. Keep the diet consistent
    and inform your doctor about eating habits.
\end{itemize}

\section{A safe shopping guide}

{\small
\begin{tabularx}{\textwidth}{|L|L|L|}
\hline
\rowcolor{verdescuro}
\textcolor{white}{\textbf{Additive/Ingredient}}
  & \textcolor{white}{\textbf{Children}}
  & \textcolor{white}{\textbf{Elderly}} \\
\hline
Azo dyes (E102, E110, E122, E124, E129) & Avoid & Limit \\
\hline
Intense sweeteners (aspartame, sucralose\ldots) & Not for young children & In moderation \\
\hline
High sodium (\textgreater1.5\,g salt/100\,g) & Limit & Avoid \\
\hline
Phosphates (E338--E452) & Moderation & Limit \\
\hline
Caffeine (energy drinks, some soft drinks) & Avoid under 12 & Caution \\
\hline
Nitrates/nitrites (E249--E252) in cured meats & Limit & Limit \\
\hline
\end{tabularx}
}

% ============================================================
% Final page with QR code
\cleardoublepage
\thispagestyle{empty}
\begin{center}
\vspace*{3cm}
{\large\bfseries\color{verdescuro} Download E-Codes Reader}\\[0.8cm]
{\normalsize Scan the QR code to find the app\\
on the Google Play Store.}\\[1.5cm]
\qrcode[height=3.5cm]{https://play.google.com/store/apps/details?id=com.faustobe.ecodes}\\[1.5cm]
\textcolor{grigioscuro}{\footnotesize\textit{faustobe.it\ \textbullet\ Version 1.0\ \textbullet\ 2026\ \textbullet\ CC BY-NC 4.0}}
\end{center}

\end{document}
